\begin{abstract}
Alternating least squares (ALS) is a standard method for training recommendation models, yet GPU implementations typically execute in FP32 on CUDA cores and leave tensor cores underutilized because rating matrices are sparse whereas tensor cores require dense low-precision tiles. We present APR-BALS, an adaptive-precision blocked ALS framework that overcomes this sparsity–tensor-core mismatch by routing entities according to their rating density. Entities accounting for the overwhelming majority of ratings are processed through tensor-core-based WMMA kernels using FP16 input operands with FP32 accumulation, while the sparse tail remains on a conventional FP32 path. A blocked batched Cholesky solver further prevents the solve stage from becoming a performance bottleneck. With a common binary, solver, and convergence criterion, APR-BALS achieves 4.1--14.8$\times$ speedup over an FP32 baseline on Netflix Prize data on an RTX 3060 as the latent dimension increases from 16 to 96, and up to 11.6$\times$ on MovieLens, while changing held-out RMSE by less than 0.005\%. Experiments on a Tesla T4 show the same performance trend. A three-way ablation demonstrates that reduced precision contributes at most 1.32$\times$ to the speedup, indicating that the primary performance gain comes from the WMMA tensor-core kernels. APR-BALS also achieves 3.0--6.4$\times$ lower per-iteration time than NVIDIA's \textit{cuMF\_als} GPU implementation on identical data splits. These results demonstrate that adaptive density-aware execution can effectively exploit tensor cores for sparse recommendation workloads while preserving FP32 accuracy.

\end{abstract}
